\documentclass[10pt]{beamer}

\usetheme{metropolis}
\definecolor{WiLabRed}{RGB}{197,18,48}
\setbeamercolor{frametitle}{fg=white,bg=WiLabRed}
\setbeamercolor{progress bar}{fg=WiLabRed!90}
\setbeamercolor{title separator}{fg=WiLabRed!90}
\setbeamercolor{progress bar in section page}{fg=WiLabRed!90}
\setbeamercolor{background canvas}{bg=white}
\setbeamercolor{alerted text}{fg=WiLabRed!90}

\usepackage{appendixnumberbeamer}

\usepackage{booktabs}
\usepackage[scale=2]{ccicons}

\usepackage{pgfplots}
\usepgfplotslibrary{dateplot}

\usepackage{xspace}
\newcommand{\themename}{\textbf{\textsc{metropolis}}\xspace}

\usepackage{marvosym}
%\usepackage{subfig}
\usepackage{graphicx}\graphicspath{{images/}}
\usepackage{subcaption}
\usepackage[framed]{./support/mcode}
\usepackage{listings}

\title{Lecture 08: Quadrature Amplitude Modulation}
\subtitle{\textit{Software Defined Radio for Engineers} (Collins~\textit{et al.}), \textsection{4.2.3}}
\date{}
\author{\textbf{Alexander M. Wyglinski, Ph.D.}}
\institute{ \vspace*{1in}\hfill\includegraphics[height=1.125cm]{wilab_logo-A70916.eps} \qquad \includegraphics[height=1.125cm]{WPI_Inst_Prim_FulClr.eps}}
% \titlegraphic{\hfill\includegraphics[height=1.5cm]{logo.pdf}}

% Foot for all slides
\setbeamertemplate{frame footer}{\tiny \copyright~2018 by Alexander Wyglinski. This work is licensed under the Creative Commons Attribution-ShareAlike 4.0 International License. To view a copy of this license, visit http://creativecommons.org/licenses/by-sa/4.0/.}

\begin{document}

%\captionsetup[subfigure]{labelformat=empty}

%%%%%%%%%%%%%%%%%%%%%%%%%%%%%%%%%%%%%%%%%%%%%%%%%%%%%%%%%%

\maketitle




%%%%%%%%%%%%%%%%%%%%%%%%%%%%%%%%%%%%%%%%%%%%%%%%%%%%%%%%%%

\frame
{
  \frametitle{Simple Receiver Structure}
  \begin{itemize}
    \item $M$-ary QAM is a popular modulation scheme due to its {\bf simple} receiver structure
    \begin{itemize}
        \item Each branch employs a $\sqrt{M}$-ary PAM detector
    \end{itemize}
  \end{itemize}

}
\frame
{
  \frametitle{Modulation Format}

    \begin{itemize}
        \item $M$-ary QAM looks like two simultaneous $\sqrt{M}$-ary PAM signal constellations
        \begin{itemize}
            \item One is acting in the real axis while the other is working on the imaginary axis
            \item For instance, a 64-QAM signal constellation can be represented by two 8-PAM signal constellations
        \end{itemize}
        \item We represent this linear modulation technique as:
        \begin{equation}
            {s_{ij}(t)}={A_i}\cdot\cos(\omega_{c}t)+{B_j}\cdot\sin(\omega_{c}t)
        \end{equation}
    \end{itemize}


}
\frame
{
  \frametitle{Computing $\varepsilon_{p,\mathsf{M-QAM}}$}

    \begin{itemize}
        \item To find the $\varepsilon_{p}$ of $M$-ary QAM, we need to determine the following:
        \begin{itemize}
            \item Calculate $d_{\min}^2=\int\limits_0^T\Delta{s^2(t)}dt=2A^2T$ using:
            \begin{equation}
            \begin{split}
                s_1(t)&=A\cdot\cos(\omega_ct)+A\cdot\sin(\omega_ct)\\
                s_2(t)&=3A\cdot\cos(\omega_ct)+A\cdot\sin(\omega_ct)\nonumber
            \end{split}
            \end{equation}
            \item For computing $\bar{E}_s$, use the expression from $M$-ary PAM and replace $M$ with $\sqrt{M}$ such that $\bar{E}_s=A^2T\frac{M-1}{3}$
            \begin{itemize}
                \item Solve for $\bar{E}_b=\frac{\bar{E}_s}{\log_2(M)}=A^2T\frac{2^b-1}{3b}$
            \end{itemize}
            \item The power efficiency is equal to $\varepsilon_{p,\mathsf{M-QAM}}=\frac{3!b}{2^b-1}$
        \end{itemize}
    \end{itemize}


}


\end{document}
